\documentclass[11pt]{article}
\usepackage[margin=1in]{geometry}
\usepackage{amsmath,amssymb,amsthm,booktabs,microtype,hyperref}
\hypersetup{hidelinks,pdfauthor={T},pdftitle={Sequential Harm Audits Under Selective Reporting}}
\newtheorem{theorem}{Theorem}
\newtheorem{proposition}{Proposition}
\title{Sequential Harm Audits Under Selective Reporting}
\author{T}
\date{}
\begin{document}
\maketitle
\begin{abstract}
Incident reports about an agent intervention may arrive more often
when the intervention causes harm. A sequential test can then raise
false alarms even when each underlying case is distinct.
We derive a sharp adjustment for a bounded ratio of reporting
probabilities and prove anytime validity with adaptive case histories.
We also give a rule for using a time-uniform upper bound learned
from calibration while preserving past report evidence. Exact
finite experiments compare the adjustments with naive monitoring
and independent audits. A paired construction shows which changes
in underlying harm remain unidentified from reports.
An independent audit top-up gives a sharp reduction in the
reporting-ratio bound when cases are deduplicated. When the full case
count is independently observed, a second test is valid even with
unbounded reporting bias. Keeping audit and report outcomes separate
also gives an exact log-optimal factor for a specified alternative.
We compare when these tests detect harm.
\end{abstract}

\section*{Summary and scope}
\paragraph{Question.} How can an operator decide when to intervene if
users report harms selectively?
\paragraph{Finding.} A conditional reporting-ratio cap gives an anytime
harm test. An upper bound learned with simultaneous calibration coverage
adds only its coverage error to the crossing bound. A pooled reporting
ratio may fail under predictable context shifts.
\paragraph{Evidence and scope.} The guarantees are proved and
checked with exact crossing calculations and observationally identical
counterexamples. The operator needs a bound for the cases on which the
monitor bets and an independent route for individual appeals.

\section{The harm rate before a case is reported}
Suppose a scheduling agent changes appointments on behalf of
its users. Some users report disruptions and others never contact
the service. The operator wants to intervene if the underlying
harm rate exceeds a reference value $p_0$, even when evidence
arrives slowly enough to make a fixed sample size impractical.
Merging duplicate reports is necessary for this question. It
still leaves the selection of which cases are reported.

Let $X_t\in\{0,1\}$ denote harm in case $t$ and
$A_t\in\{0,1\}$ indicate whether its outcome enters the audit.
The observed outcome is either omitted, admitted clean, or admitted
harmful. The history $\mathcal F_{t-1}$ includes information available
when the monitor sets its next betting factor. We consider the
conditional null $\Pr(X_t=1\mid\mathcal F_{t-1})\le p_0$.
The selection model permits different admission probabilities
for harmful and clean cases, with ratio at most $\rho$ at each history.

The main result separates two choices an operator can make.
The reference rate $p_0$ expresses a tolerance for underlying harm.
The reporting ratio $\rho$ reflects the audit's observation process.
This response probability ratio is a sensitivity parameter for
outcome-dependent observation. Smith and VanderWeele \cite{smith}
bound causal risk-ratio distortion under selection. Here we use
a conditional reporting ratio to construct a sequential test,
with a nonnegative supermartingale as in time-uniform inference
\cite{howard}.

\section{A sharp conditional reporting bound}
\begin{theorem}
Let $0<p_0<1$ and $1\le\rho<\infty$. At every history let
$a_x=\Pr(A_t=1\mid X_t=x,\mathcal F_{t-1})$, with
$a_1\le\rho a_0$. Set
\[
q=\frac{\rho p_0}{1-p_0+\rho p_0}.
\]
When $\Pr(A_t=1\mid\mathcal F_{t-1})>0$, the admitted
harm probability is at most $q$ under the conditional null.
Let $b_t$ be any predictable value in $[q,1]$. Start with $M_0=1$
and multiply by $1$ on omission, by $b_t/q$ on admitted harm,
and by $(1-b_t)/(1-q)$ on admitted clean outcomes.
For any $0<\alpha<1$,
\[
\Pr\!\left(\exists t\ge0:\ M_t\ge 1/\alpha\right)\le\alpha .
\]
The bound $q$ on admitted harm is sharp.
\end{theorem}
\begin{proof}
Fix a history and write $p=\Pr(X_t=1\mid\mathcal F_{t-1})$.
At histories with positive admission, Bayes' rule gives
\[
\Pr(X_t=1\mid A_t=1,\mathcal F_{t-1})
 =\frac{p a_1}{p a_1+(1-p)a_0}
 \le \frac{\rho p}{1-p+\rho p}
 \le q .
\]
The odds of harm have been multiplied by at most $\rho$.
Positive admission under the null implies $a_0>0$, since
$a_0=0$ forces $a_1=0$ by assumption. Null conditional events
can be assigned arbitrary consistent versions of $a_x$.

Write $r$ for the admitted harm probability. The mean factor
on admission is
\[
r\,\frac{b_t}{q}+(1-r)\frac{1-b_t}{1-q}
 =1+(r-q)\frac{b_t-q}{q(1-q)}\le1 .
\]
Omissions contribute factor one, so $M_t$ is a nonnegative
supermartingale. Stop it at its first crossing or a fixed finite
horizon. Its expected stopped value is at most one, while on
crossing its value is at least $1/\alpha$. Letting the horizon
increase gives the result. For sharpness, take $p=p_0$,
$a_1=1$ and $a_0=1/\rho$.
\end{proof}

The factor can use a history-dependent bet without changing the
argument. A large $\rho$ brings $q$ close to one and makes
successive harmful reports weak evidence. A process that learns
how strongly to bet may improve power, but it still needs an
externally justified upper bound on selection. The rate $q$ can
also vary predictably with the case class when the reference
rate and reporting bound are specified conditionally.

\section{Reducing selection bias with independent audit draws}
Let $R$ indicate an adjudicated spontaneous report. Before using
any current-case evidence, choose an audit probability $s$ from
the past history. Select the audit with probability $s$ independently
of both the harm and the report, conditional on that history.
Admit a case if either route provides the outcome. If both routes
apply, record it once. An audit assignment can reuse an already
adjudicated outcome. This requires comparable outcome definitions
on both routes.

\begin{proposition}
If spontaneous reporting satisfies $a_1\le\rho a_0$, the
combined admission probabilities satisfy the sharp ratio bound
\[
\rho_{\rm pool}=\frac{\rho}{1+s(\rho-1)}.
\]
With no finite spontaneous reporting bound and $s>0$, the
sharp combined ratio bound is $1/s$.
\end{proposition}
\begin{proof}
The combined probability for outcome $x$ is $s+(1-s)a_x$.
For fixed $a_0$, the largest numerator uses
$a_1=\min(1,\rho a_0)$. Up to $a_0=1/\rho$, the ratio is
nondecreasing in $a_0$. Beyond that point the numerator is one,
and the ratio is nonincreasing. The maximum is attained at
$(a_0,a_1)=(1/\rho,1)$. With unrestricted reporting the
largest numerator is one and the smallest denominator is $s$,
attained by $(0,1)$.
\end{proof}

Use $\rho_{\rm pool}$ in the conditional harm bound at each
history. For $\rho>1$, a target combined ratio $r\in[1,\rho]$
requires $s\ge(\rho/r-1)/(\rho-1)$. The expected number
of additional adjudications is $s\Pr(R=0)$ per opportunity if
already reported cases have a usable label. A draw concentrated
on already suspected harms needs a different observation model.
Combining two routes for one case does not create two independent
outcomes.

The verifier compared 4,514 grid ratios, with a sharp case
for each of 36 settings of $\rho$ and $s$. For the existing
ratio-three example, set $s=1/4$, giving a combined ratio of
two. Over 48 opportunities with threshold eight, use the same
harm alternative $4/5$ to set a fixed bet for each observation
rule. The resulting crossing probabilities are
\begin{center}
\begin{tabular}{lrr}
\toprule
Observation rule & Null at $1/2$ & Alternative at $4/5$ \\
\midrule
Spontaneous reports & 0.0888 & 0.9013 \\
Deduplicated pool & 0.0958 & 0.9543 \\
Independent assignments only & 0.0606 & 0.6469 \\
\bottomrule
\end{tabular}
\end{center}
The pooled and assignment-only rules use the same random
assignments. The pool expects four extra adjudications on
unreported cases under the null and $8/5$ under the alternative.
These are model-based costs and crossing probabilities for
the stipulated reporting laws.

\section{Using all enumerated opportunities}
Suppose the monitor observes every eligible case even when neither
route provides a label. Let $Y=1$ mean that an adjudicated harmful
outcome entered the audit and let $Y=0$ for all other cases. Since
$Y\le X$, the conditional null also bounds $\Pr(Y=1)$ by $p_0$.
This gives a valid betting factor for any predictable $d\in[p_0,1]$,
using $d/p_0$ on $Y=1$ and $(1-d)/(1-p_0)$ otherwise.
It requires a complete case count and no false harmful labels.
It uses omissions as evidence about the detected harm rate.
If reporting suppresses so many harms that this rate remains
below $p_0$, it cannot provide the intended power.

The monitor can also retain the route by which it obtained each
label. With audit probability $s$ and spontaneous reporting
probabilities $r_0,r_1$, there are five outcomes. They are audited
harm, audited clean, unaudited reported harm, unaudited reported
clean, and omission. Their probabilities are
\[
 \bigl(sp,\ s(1-p),\ (1-s)pr_1,\ (1-s)(1-p)r_0,\
 (1-s)[p(1-r_1)+(1-p)(1-r_0)]\bigr).
\]
An independently assigned audit supplies a usable label in this
model. Count each case on exactly one route.

\begin{proposition}
For known $s$ and finite $\rho\ge1$, a nonnegative five-outcome
factor has conditional expectation at most one for every null law
with $p\le p_0$ and $r_1\le\rho r_0$ exactly when it passes that
check at $p\in\{0,p_0\}$ and
\[
 (r_0,r_1)\in\{(0,0),(1,0),(1,1),(1/\rho,1)\}.
\]
\end{proposition}
\begin{proof}
For fixed reporting probabilities the expectation is affine in
$p$, so its maximum is at one endpoint. At each endpoint it is
linear on the reporting polygon $0\le r_0,r_1\le1$,
$r_1\le\rho r_0$. The listed points are the vertices, with a
repeat when $\rho=1$. Every vertex law is allowed by the null.
\end{proof}

For example, use the ordinary harm threshold in the assigned
branch and the report-adjusted threshold in the unassigned branch.
The five-outcome check certifies this rule and allows other factors
that use the route information. The verifier checked 1,152
conditional expectations for denominator, pooled and route-specific
factors at the extreme null laws.

For the denominator and pooled rules, a fixed weighted average
of their wealth processes is also a nonnegative supermartingale.
The same anytime threshold applies to that average. It shares
evidence across the two rules without treating their outcomes as
independent. The following exact comparison uses equal weights,
24 opportunities, threshold eight, $s=1/4$, and $\rho=3$.
The denominator bet is $d=4/5$ and the pooled bet is $8/9$.
The null uses $p=1/2$ with $(r_0,r_1)=(1/3,1)$. Both
alternatives use $p=4/5$, with reporting pairs $(1/3,1)$
and $(1/10,3/10)$ respectively.
\begin{center}
\begin{tabular}{lrrr}
\toprule
Case law & Denominator & Pooled & Equal mixture \\
\midrule
Null & 0.0826 & 0.0730 & 0.0788 \\
All harms reported & 0.8978 & 0.7945 & 0.8999 \\
Sparse reports & 0.0144 & 0.2970 & 0.1545 \\
\bottomrule
\end{tabular}
\end{center}
The mixture gives up power relative to the pooled rule in the
sparse fixture. A service can choose the mix before evaluating
the live stream, using separate evidence about how complete its
case and reporting records are. Exact path enumeration agreed
with the crossing recursion on nine short runs.

\section{Exact finite comparison}
We computed first-crossing probabilities for 48 opportunities,
threshold eight, and independent opportunities within each fixture.
The naive process uses $q=1/2$ and $b=3/4$. The adjusted process
uses $q=3/4$ and $b=7/8$. The independent audit selects each case
with probability $2/3$ and applies the $1/2$ reference. In the
primary null the expected sample size matches that of admitting
all harmful cases and one third of clean cases.

\begin{center}
\begin{tabular}{llrr}
\toprule
Selection & Monitor & $p=1/2$ & $p=4/5$ \\
\midrule
Ratio $3$ & Naive & 0.8809 & 1.0000 \\
Ratio $3$ & Adjusted & 0.0620 & 0.9188 \\
Ratio $4$ & Adjusted using $3$ & 0.1528 & 0.9672 \\
Independent audit & $1/2$ reference & 0.0919 & 0.9659 \\
\bottomrule
\end{tabular}
\end{center}
Table entries are rounded from exact rationals. The algorithm
propagates surviving probability mass over pairs $(h,c)$ of admitted
harmful and clean counts. On each transition it removes mass when
$(b/q)^h((1-b)/(1-q))^c\ge8$. The omitted transition leaves the
pair unchanged. The verifier checks conservation of mass and
compares a seven-step case with independent enumeration of all
$3^7$ paths. The ratio-four condition intentionally violates the
assumed bound. Its null crossing probability exceeds $1/8$ in
this finite experiment. Under the specified alternative and the
primary selection model, the corrected monitor crosses in about
$92\%$ of cases. This power belongs to that fixture.

\section{An identification limit}
\begin{proposition}
A known upper bound $\rho=3$ on the response probability ratio
does not generally identify whether the underlying harm rate
exceeds $1/2$. This remains true when harms are at least as likely
to be reported as clean outcomes.
\end{proposition}
\begin{proof}
In the null world take $(p,a_1,a_0)=(1/2,3/5,3/10)$.
In the alternative take $(3/5,1/2,3/8)$. Each world has an
admitted harm with probability $3/10$, an admitted clean
outcome with probability $3/20$, and an omission with probability
$11/20$. Both satisfy $1\le a_1/a_0\le3$.
Repeating each world independently gives identical laws for all
observed histories. Consequently any test based on those histories
has identical crossing probabilities in the two worlds.
\end{proof}

\section{Learning an upper bound without restarting the monitor}
Suppose a fixed $\rho_\star$ bounds the conditional reporting
ratio at all histories, but its value is unknown. Write
$q(\rho)=\rho p_0/(1-p_0+\rho p_0)$ and $q(+\infty)=1$.
Use a predictable coefficient $\lambda_t\in[0,1]$, chosen
once at each time for all candidate ratios. For a candidate
$\rho$, let $E_t(\rho)$ be the product of factors
\[
 1\quad\text{on omission},\qquad
 1+\lambda_t(X_t-q(\rho))\quad\text{on admission}.
\]
These factors are nonnegative and decrease with $\rho$
on every realized path. For $\rho_\star$, the conditional
mean factor on admission is at most one, so
$E_t(\rho_\star)$ is a nonnegative supermartingale under the
null. Keep the historical coefficients fixed when evaluating a
new candidate ratio.

\begin{theorem}
Let calibration produce upper bounds $U_t\ge1$ with
$\Pr(\exists t:U_t<\rho_\star)\le\beta$. Then stopping when
$E_t(U_t)\ge1/\alpha$ has null crossing probability at most
$\alpha+\beta$.
\end{theorem}
\begin{proof}
On the event of simultaneous calibration coverage, pathwise
monotonicity gives $E_t(U_t)\le E_t(\rho_\star)$ for all $t$.
On that event, crossing by the calibrated process implies
crossing by the fixed-ratio process, which has probability at
most $\alpha$. Add the probability at most $\beta$ that the
upper bound ever fails. The argument needs no independence
between the two processes when the conditional null and
ratio bound hold in the monitor's full pre-case filtration.
The guarantee comes from comparison with $E_t(\rho_\star)$.
The plug-in process itself need not be a supermartingale.
\end{proof}

For example, with fixed reporting probabilities $a_0,a_1$,
representative independent audit samples within each outcome
stratum can give a simultaneous upper bound. After $n$ audited
cases in a stratum, use one-sided Hoeffding radius
\[
 d_n=\sqrt{\frac{\log(2n(n+1)/\beta)}{2n}}.
\]
Take the lower endpoint for $a_0$ and the upper endpoint for
$a_1$, clipped to the unit interval. Their respective error
probabilities are each at most $\beta/(2n(n+1))$. Summing over
both strata and every $n$ gives total probability at most
$\beta$. Divide the upper endpoint by the lower one, flooring
the result at one and using $+\infty$ when the lower endpoint
is zero or a stratum is missing. The resulting upper bound
can be reevaluated as audits arrive. A narrower interval does
not justify changing the coefficients of earlier bets.

\paragraph{A pooled calibration can fail.}
Let two observed contexts be equally common, with harm
probability $1/2$ in both. Their $(a_0,a_1)$ pairs are
$(1/4,1)$ and $(3/4,0)$. Pooled reporting probabilities are
equal, yet the first context has ratio four. A monitor that
bets only in that context, uses ratio one, and sets
$\lambda_t=1$ has conditional mean factor $19/16$ under
the null there. An upper bound on a pooled reporting ratio
cannot replace a bound at histories where the monitor bets.

\section{An exact calibration experiment}
In a finite two-point model, the clean reporting probability
is $1/4$ and the true ratio is one or three. Before opportunity
four, use $U_t=3$. At that opportunity reveal the reporting
indicators from two independently adjudicated harmful cases.
Lower $U_t$ to one for the remaining decisions only if both
were unreported. At true ratio three, this happens with
probability $1/16$, which bounds simultaneous undercoverage
in this experiment. Calibration therefore spends two extra
known-harm audits. We use coefficient one and threshold eight
over 40 opportunities. The theorem's error allowance is
$1/8+1/16=3/16$.

\begin{center}
\begin{tabular}{lrrrr}
\toprule
True setting & Fixed upper three & Calibrated & Known ratio & Naive ratio one\\
\midrule
Null, ratio one & 0.0005 & 0.0232 & 0.0408 & 0.0408\\
Null, ratio three & 0.0797 & 0.1205 & 0.0797 & 0.7324\\
Alternative, ratio one & 0.0608 & 0.3241 & 0.5289 & 0.5289\\
Alternative, ratio three & 0.7605 & 0.7754 & 0.7605 & 0.9997\\
\bottomrule
\end{tabular}
\end{center}
The alternative is $p=4/5$. Entries are rounded exact
first-crossing probabilities. At true ratio one, calibration
recovers part of the power lost by always using three. At
ratio three, occasional undercoverage raises the null crossing
probability relative to knowing the true ratio. These operating
points belong to the declared two-point model. The calculation
agrees with 8,748 six-step path evaluations across twelve
oracle settings. The full receipt retains the exact values.

\section{A factor tuned to sparse reporting}
A monitor that retains the observation route can target a specified
alternative. Let $Q$ be its strictly positive five-outcome law, and
let $P^*$ be a strictly positive mixture of admissible null laws.
Consider the candidate factor $e^*(i)=Q_i/P_i^*$. If its mean is
at most one at each null vertex, it is valid and has the largest
expected log factor under $Q$ among all valid factors. This is
the finite form of the reverse information projection construction
of Larsson, Ramdas, and Ruf \cite{larsson}. Indeed, any valid $e$
has $\sum_i P_i^*e(i)\le1$, so Jensen's inequality gives
\[
 E_Q[\log(e/e^*)]\le \log\!\left(\sum_i P_i^*e(i)\right)\le0.
\]

Here is a family where the factor is explicit. Fix $p_0<p_1<1$ and
an audit probability $0<s<1$. In the alternative set $r_1=\rho r_0$
and let $m$ be its total spontaneous reporting probability. Assume
$0<m<1$ and $m\le p_0+(1-p_0)/\rho$. For the null candidate
use $p_0$ with the same $m$ and the same reporting ratio. Both laws
are feasible because their harm-report probability is obtained from
$r_1=\rho m/(1-p+\rho p)\le1$. Set
$q_j=\rho p_j/(1-p_j+\rho p_j)$. The tagged factor is
\[
 \left(\frac{p_1}{p_0},\ \frac{1-p_1}{1-p_0},\
       \frac{q_1}{q_0},\ \frac{1-q_1}{1-q_0},\ 1\right).
\]
Under the null, the harm probability is at most $p_0$ on the
audit route and at most $q_0$ conditional on an unaudited report.
Both conditional mean factors are at most one. Omissions have
factor one. This proves validity and, by the preceding argument,
log optimality for the stated alternative.

For the earlier sparse fixture the factor is
$(8/5,2/5,16/13,4/13,1)$. The 24-opportunity crossing
probability is $0.4956$, compared with $0.2970$ for the pooled
factor and $0.0144$ for the denominator factor. In the dense
alternative, using the sparse factor lowers the crossing probability
to $0.7634$ from the denominator's $0.8978$. For that dense
alternative the denominator factor is itself log optimal, by
taking $P^*$ with the dense reporting parameters and $p=p_0$.
The sparse factor crossed with probability $0.0746$ and $0.0514$
in the dense and sparse null fixtures. All use the same threshold
of eight. The verifier checked 304 vertex constraints and reproduced
12 short crossings by full path enumeration. Log optimality
concerns expected growth under the declared alternative. The
finite-horizon comparisons concern the individual simulated laws.

A monitor can use past or separate evidence to choose a factor
before the current case. Fixing the factor for sparse reports
has a cost when reporting is dense. Retaining route labels gives
a choice between these factors while preserving the null test
as long as each chosen factor satisfies the conditional constraint.

\section{Using the audit to correct an agent}
An operator can use independent sampling to investigate
silent harm or use a justified response ratio to interpret existing
reports. The first choice costs review time. The second can reduce
that cost where the reporting channel is well characterized.
A separate appeal route remains necessary because a statistical
alarm identifies a pattern of outcomes rather than adjudicating
a specific person's case. A practical first deployment would be
a shadow audit on paired notification or scheduling interventions,
with an independent random sample used to estimate reporting
differences. Stratify by how easily users can report. The
prospective test should compare alarm timing, missed harms, and
review cost while preserving an independently observed denominator.
Unverified case identities, delayed outcome adjudication, and
selection that shifts after an alarm require additional controls.

\paragraph{Review status.} Independent manuscript review is pending.

\begin{thebibliography}{9}
\bibitem{smith}
Louisa H. Smith and Tyler J. VanderWeele.
\emph{Bounding Bias Due to Selection}. Epidemiology 30(4),
509--516, 2019. \url{https://arxiv.org/abs/1810.13402}.
\bibitem{howard}
Steven R. Howard, Aaditya Ramdas, Jon McAuliffe, and Jasjeet Sekhon.
\emph{Time-uniform, nonparametric, nonasymptotic confidence sequences}.
2021. \url{https://arxiv.org/abs/1810.08240}.
\bibitem{larsson}
Martin Larsson, Aaditya Ramdas, and Johannes Ruf.
\emph{The numeraire e-variable and reverse information projection}.
arXiv:2402.18810, version 4, 2025.
\url{https://arxiv.org/pdf/2402.18810}.
\end{thebibliography}
\end{document}
